\chapterstar{Conclusiones}
El desarrollo e implementación de la plataforma cloud de gestión de inventario y monitoreo ambiental para Farmacia Selene culminó con el cumplimiento íntegro del objetivo general y de los objetivos específicos planteados. Se logró diseñar una solución tecnológica que unifica la administración operativa y la supervisión física del entorno, resolviendo la desorganización y la falta de información actualizada que caracterizaban a los procesos manuales previos.  La base de datos estructurada en PostgreSQL cumplió el objetivo de centralizar la información de productos, lotes, ventas y parámetros ambientales.  Su diseño modular garantizó la integridad referencial y permitió implementar reglas de negocio complejas, como el uso de índices parciales para gestionar medicamentos sin código de lote asignado.  El servidor desarrollado con Python y FastAPI procesó eficientemente las solicitudes del sistema y fue desplegado con éxito en la plataforma Railway.  La lógica del servidor automatizó exitosamente la política FEFO (Primeras Entradas, Primeras Salidas), asegurando que el sistema descuente el inventario priorizando los lotes con la caducidad más próxima.  La interfaz de usuario, desarrollada con React y Material-UI, cumplió el objetivo de proporcionar un entorno visual accesible y responsivo desplegado en Vercel.  Esta capa de la aplicación permitió al personal interactuar de forma fluida con los módulos de inventario, ventas y el tablero de monitoreo sin necesidad de instalar software local.  La capa física, gobernada por el microcontrolador ESP32-C5, cumplió su función de capturar y transmitir datos ambientales al servidor mediante solicitudes HTTP cada 30 segundos.  El sensor AHT10 demostró ser preciso para validar las condiciones del entorno y generar alertas cuando los valores superaron los umbrales de la normatividad.  La integración del lector de códigos de barras se resolvió exitosamente mediante la implementación de un "modo lector" activado con la tecla F2.  Este desarrollo agilizó el registro de productos durante las transacciones y minimizó el margen de error humano en el punto de venta.  A pesar de los retos técnicos encontrados durante la integración física, como el conflicto de pines con el LED RGB y las limitaciones de altura del conector Dupont en el gabinete impreso en 3D, los problemas fueron diagnosticados y resueltos o se plantearon rutas de rediseño efectivas.  

\thispagestyle{empty}


\chapterstar{Recomendaciones}


\thispagestyle{empty}


\chapterstar{Mejoras}


\thispagestyle{empty}